The half-chord is eight, so its distance from the origin is $\sqrt{17^2-8^2}=15$. The common chord line can be $x=15$ or $x=-15$. Its distance from the second centre $(21,0)$ is respectively six or 36, giving radii $\sqrt{6^2+8^2}=10$ or $\sqrt{36^2+8^2}=4\sqrt{85}$. Thus the source allows $\boxed{10\text{ or }4\sqrt{85}}$. If the common chord is stipulated to lie between the centres, only ten remains.
\begin{asy}
import graph;import geometry;usepackage("amsmath");size(7cm);
draw(circle((0,0),17));draw(circle((21,0),10),blue);draw((15,-8)--(15,8),red);dot((0,0));dot((21,0));label("$O$",(0,0),(0,-1));label("$(21,0)$",(21,0),(0,-1));label("between-centres case",(5,-22));
\end{asy}

\begin{asy}
import graph;import geometry;usepackage("amsmath");size(7cm);
pair O=(0,0),P=(21,0);draw(circle(O,17));draw(circle(P,4*sqrt(85)),blue);draw((-15,-8)--(-15,8),red);dot(O);dot(P);label("$O$",O,(0,-1));label("$P$",P,(0,-1));label("common chord outside the centre segment",(21,-42));
\end{asy}

